% Options for packages loaded elsewhere
% Options for packages loaded elsewhere
\PassOptionsToPackage{unicode}{hyperref}
\PassOptionsToPackage{hyphens}{url}
\PassOptionsToPackage{dvipsnames,svgnames,x11names}{xcolor}
%
\documentclass[
  letterpaper,
  DIV=11,
  numbers=noendperiod]{scrartcl}
\usepackage{xcolor}
\usepackage{amsmath,amssymb}
\setcounter{secnumdepth}{-\maxdimen} % remove section numbering
\usepackage{iftex}
\ifPDFTeX
  \usepackage[T1]{fontenc}
  \usepackage[utf8]{inputenc}
  \usepackage{textcomp} % provide euro and other symbols
\else % if luatex or xetex
  \usepackage{unicode-math} % this also loads fontspec
  \defaultfontfeatures{Scale=MatchLowercase}
  \defaultfontfeatures[\rmfamily]{Ligatures=TeX,Scale=1}
\fi
\usepackage{lmodern}
\ifPDFTeX\else
  % xetex/luatex font selection
\fi
% Use upquote if available, for straight quotes in verbatim environments
\IfFileExists{upquote.sty}{\usepackage{upquote}}{}
\IfFileExists{microtype.sty}{% use microtype if available
  \usepackage[]{microtype}
  \UseMicrotypeSet[protrusion]{basicmath} % disable protrusion for tt fonts
}{}
\makeatletter
\@ifundefined{KOMAClassName}{% if non-KOMA class
  \IfFileExists{parskip.sty}{%
    \usepackage{parskip}
  }{% else
    \setlength{\parindent}{0pt}
    \setlength{\parskip}{6pt plus 2pt minus 1pt}}
}{% if KOMA class
  \KOMAoptions{parskip=half}}
\makeatother
% Make \paragraph and \subparagraph free-standing
\makeatletter
\ifx\paragraph\undefined\else
  \let\oldparagraph\paragraph
  \renewcommand{\paragraph}{
    \@ifstar
      \xxxParagraphStar
      \xxxParagraphNoStar
  }
  \newcommand{\xxxParagraphStar}[1]{\oldparagraph*{#1}\mbox{}}
  \newcommand{\xxxParagraphNoStar}[1]{\oldparagraph{#1}\mbox{}}
\fi
\ifx\subparagraph\undefined\else
  \let\oldsubparagraph\subparagraph
  \renewcommand{\subparagraph}{
    \@ifstar
      \xxxSubParagraphStar
      \xxxSubParagraphNoStar
  }
  \newcommand{\xxxSubParagraphStar}[1]{\oldsubparagraph*{#1}\mbox{}}
  \newcommand{\xxxSubParagraphNoStar}[1]{\oldsubparagraph{#1}\mbox{}}
\fi
\makeatother


\usepackage{longtable,booktabs,array}
\newcounter{none} % for unnumbered tables
\usepackage{calc} % for calculating minipage widths
% Correct order of tables after \paragraph or \subparagraph
\usepackage{etoolbox}
\makeatletter
\patchcmd\longtable{\par}{\if@noskipsec\mbox{}\fi\par}{}{}
\makeatother
% Allow footnotes in longtable head/foot
\IfFileExists{footnotehyper.sty}{\usepackage{footnotehyper}}{\usepackage{footnote}}
\makesavenoteenv{longtable}
\usepackage{graphicx}
\makeatletter
\newsavebox\pandoc@box
\newcommand*\pandocbounded[1]{% scales image to fit in text height/width
  \sbox\pandoc@box{#1}%
  \Gscale@div\@tempa{\textheight}{\dimexpr\ht\pandoc@box+\dp\pandoc@box\relax}%
  \Gscale@div\@tempb{\linewidth}{\wd\pandoc@box}%
  \ifdim\@tempb\p@<\@tempa\p@\let\@tempa\@tempb\fi% select the smaller of both
  \ifdim\@tempa\p@<\p@\scalebox{\@tempa}{\usebox\pandoc@box}%
  \else\usebox{\pandoc@box}%
  \fi%
}
% Set default figure placement to htbp
\def\fps@figure{htbp}
\makeatother





\setlength{\emergencystretch}{3em} % prevent overfull lines

\providecommand{\tightlist}{%
  \setlength{\itemsep}{0pt}\setlength{\parskip}{0pt}}



 


\KOMAoption{captions}{tableheading}
\makeatletter
\@ifpackageloaded{caption}{}{\usepackage{caption}}
\AtBeginDocument{%
\ifdefined\contentsname
  \renewcommand*\contentsname{Table of contents}
\else
  \newcommand\contentsname{Table of contents}
\fi
\ifdefined\listfigurename
  \renewcommand*\listfigurename{List of Figures}
\else
  \newcommand\listfigurename{List of Figures}
\fi
\ifdefined\listtablename
  \renewcommand*\listtablename{List of Tables}
\else
  \newcommand\listtablename{List of Tables}
\fi
\ifdefined\figurename
  \renewcommand*\figurename{Figure}
\else
  \newcommand\figurename{Figure}
\fi
\ifdefined\tablename
  \renewcommand*\tablename{Table}
\else
  \newcommand\tablename{Table}
\fi
}
\@ifpackageloaded{float}{}{\usepackage{float}}
\floatstyle{ruled}
\@ifundefined{c@chapter}{\newfloat{codelisting}{h}{lop}}{\newfloat{codelisting}{h}{lop}[chapter]}
\floatname{codelisting}{Listing}
\newcommand*\listoflistings{\listof{codelisting}{List of Listings}}
\makeatother
\makeatletter
\makeatother
\makeatletter
\@ifpackageloaded{caption}{}{\usepackage{caption}}
\@ifpackageloaded{subcaption}{}{\usepackage{subcaption}}
\makeatother
\usepackage{bookmark}
\IfFileExists{xurl.sty}{\usepackage{xurl}}{} % add URL line breaks if available
\urlstyle{same}
% fallback for those not using the hyperref driver hyperxmp:
\makeatletter
\@ifundefined{xmpquote}{\newcommand{\xmpquote}[1]{#1}}{}
\makeatother
\hypersetup{
  pdftitle={My Quarto Document},
  colorlinks=true,
  linkcolor={blue},
  filecolor={Maroon},
  citecolor={Blue},
  urlcolor={Blue},
  pdfcreator={LaTeX via pandoc}}


\title{My Quarto Document}
\author{}
\date{}
\begin{document}
\maketitle


Teaching Philosophy

I am a recent PhD in philosophy from Saint Louis University. My work has
been in philosophy of mind, cognitive science, and moral psychology
where I focus on questions regarding the formation of human agency and
personhood. I defended my PhD dissertation on moral epistemology and
narratives. I attribute my philosophical pedagogy to my dissertation
advisor Helen De Cruz, and her work on fiction and philosophy which
played an insurmountable role in my philosophical training. In addition
my co-advisor Dan Haybron who's work in value theory and happiness was a
large influence on my teaching and research, as well as Eleonore Stump
whose work on Thomistic virtues have been a large influence on my
understanding of the importance of narrative. In my academic work, I
focus broadly on narratives themselves, and what narratives can do for
human moral knowledge and understanding. I have taught undergraduate
courses in Ethics, Critical Thinking, introduction to philosophy, and
Self and Reality, philosophy of law, aesthetics, and philosophy of mind.
In the spring of 2027 I will be teaching my first ever course in
personhood and identity where I will focus on an account of social
identity grounded in narrative knowledge of one's self. As part of a
larger project tying together sentiment analysis and moral foundations
theories, I have also taught graduate level workshops in sentiment
analysis, R and quanteda. This project attempted to tie together social
epistemology and Aristotelian virtue ethics through a focus on public
narratives. I have also designed and delivered several custom
asynchronous Philosophy-Digital Humanities courses, where I integrated
teaching methods relying on technology to enhance student engagement and
critical thinking skills to about 80 students. This summer, I taught an
asynchronous introduction to philosophy course to 10 students. This
course was based on David Chalmers' Reality+ and in it, I used the open
world video game Elite Dangerous to explore Dave Chalmers' and Andy
Clark's extended cognition thesis, as well as explore concepts of
meaning in virtual environments. I believe that for a concept to be
memorable, it must hold intrinsic value for the student. In my courses,
I leverage this fact, providing opportunities for my students to learn
from the material that I carefully and thoughtfully curate, and
conversations about the material. As such, I push my students to reflect
on storytelling and personal engagement. In my student evaluations, my
students have noted that they often find the course material and
discussions with each other memorable. In evaluations, students have
often reflected that my approach helps them to see philosophical
questions as central to a meaningful life. In my classrooms, I pursue
the following values in my courses. Critical thinking, listening,
self-inquiry and reflection. These are central to understanding. Based
on my experiences in the digital humanities, I use artifacts from
popular culture, songs, plays, video games, etc., in my assignments to
motivate students to apply their experiences outside of the classroom to
the course material. Which promotes personal engagement with course
materials that in turn prompts the student to pursue philosophical
insight for its own benefit. This approach is interdisciplinary, and
blends philosophical concepts with digital tools, and focuses on a
dynamic learning environment for diverse student groups. In addition, I
have implemented adaptive learning strategies such as D\&D campaigns,
tailoring course content to individual student needs, resulting in mixed
results. and facilitate online discussions and group projects by
assigning projects that required the student to create a philosophical
dialogue between two or more fictional characters exploring some
concept. The intention was that through group assignments, the students
would create a collaborative virtual classroom that promotes
peer-to-peer learning and knowledge-sharing. I also work towards
developing a forward-thinking curriculum that bridges traditional
philosophical inquiry with contemporary digital literacy. This prepares
students for an evolving academic landscape that requires an increasing
familiarity with digital and computational technologies. With the help
of members of MIHLab.org, I was able to create a turn-based dialogue
game in Unity. I introduced this game to my students as a study tool. In
it, the students would enter the video game where they would be
confronted with questions and prompts from fictional figures based on
our readings. My most recent pedagogical projects look to incorporate
game play footage and game play as a form of philosophical thought
experiment into course discussions. I am currently a member of the Saint
Louis Digital Humanities Network. This is a group of educators in the
Saint Louis area that are committed to exploring ways and implications
of expanding use of technology for pedagogical purposes. Given that my
philosophical interests are motivated by questions about our social
environments, these new modes of social interaction are of great
interest. With this in mind, I created the Moral Imagination and Hope
Laboratory. At MIHLab, we design philosophy thought experiments using
Mass Multiplayer Online / Open World video-games (Pilot, fall 2025) as
well as create interactive technologies such as video games and
curricula. I hope to collectively explore implications for conducting
our lives in digitally represented social environments by using
philosophical reflection to better understand how online communities can
be used as a social learning environment. It is our hope that we can use
emerging technologies to create lasting communities of critical thinkers
and civic-minded future generations. I am excited about the opportunity
to contribute to a learning environment that seeks to engage the student
population in meaningful philosophical discourse and prepare them for a
changing world. My interdisciplinary experience, and philosophical rigor
combined with innovative pedagogical methods will be a terrific addition
to your department. I would welcome the opportunity to discuss my
qualifications and the needs of your department further. Thank you for
your time, and for your thoughtful consideration.




\end{document}
